%!TEX TS-program = pdflatexmk
%the line above is only for offline compilation (assuming your system is set up for it)
%this will be based on the amsart class, which probably causes more pain than necessary
%but gets a lot of the formatting the way I like it
%look for !!!EDIT comments
\documentclass[oneside,12pt]{amsart}

%format
\usepackage[left=1in,top=1in,right=1in,bottom=1in,footskip = 0.333in]{geometry} %add showframe option to show true border
\usepackage[T1]{fontenc}
\usepackage[utf8]{inputenc}

%math
\usepackage{nicefrac} %provides nicefrac env
\usepackage{amsmath,amsfonts,amsthm,amssymb}

%utility
\usepackage{longtable} %provides longtable env
\usepackage[hyphens]{url} %include urls but allow breaks at hyphens (for footnotes)
\usepackage{tabularx} %provides tabularx env
\newcolumntype{L}[1]{>{\raggedright\let\newline\\\arraybackslash\hspace{0pt}}m{#1}} %define left justified column type

%figures
\usepackage{graphicx}
\usepackage[margin=0.1cm,textfont=it]{caption}
\usepackage{subcaption}
\captionsetup[subfigure]{labelfont=rm}
\usepackage{wrapfig}

%create \nopar command to force a paragraph break without extra spacing
%useful when trying to get wrapfigure to behave
\newcommand{\nopar}{{\parfillskip=0pt \parskip=0pt \par}}
%useful command to issue inside math env to play nice with wrapfig:
%\displaywidth=\parshapelength\numexpr\prevgraf+2\relax

%include pdfs
% \usepackage{pdflscape} %only needed if doing a landscape PDF
% \usepackage{rotating}
\usepackage{pdfpages}

%%for Draft watermark (editing phase only)
%\usepackage{draftwatermark}
%\SetWatermarkText{DRAFT}
%\SetWatermarkScale{5}

%set up header (this has to come before the section patching)
\usepackage{fancyhdr}
\setlength{\headheight}{0.2in}
\pagestyle{fancy}
\fancyhf{}
\lhead{\small \sc Program} %!!!EDIT
\chead{\small \sc \rightmark}
\rhead{\footnotesize \thepage}
\renewcommand{\headrulewidth}{0pt}

% need to redefine section to allow rightmark to work with amsart
\let\origsection\section
\renewcommand{\section}[1]{\sectionmark{#1}\origsection{#1}}
% redefine sectionmark
\renewcommand\sectionmark[1]{\markboth{#1}{#1}}

%fix toc
\setcounter{tocdepth}{4}% to get subsubsections in toc use 3, 4 to include paragraphs too
\makeatletter
\renewcommand{\@pnumwidth}{2em}% default is 1.55em -accommodates wider page numbers
\def\@tocline#1#2#3#4#5#6#7{\relax
  \ifnum #1>\c@tocdepth % then omit
  \else
    \par \addpenalty\@secpenalty\addvspace{#2}%
    \begingroup \hyphenpenalty\@M
    \@ifempty{#4}{%
      \@tempdima\csname r@tocindent\number#1\endcsname\relax
    }{%
      \@tempdima#4\relax
    }%
    \parindent\z@ \leftskip#3\relax \advance\leftskip\@tempdima\relax
    \rightskip\@pnumwidth plus4em \parfillskip-\@pnumwidth
    #5\leavevmode\hskip-\@tempdima
      \ifcase #1
       \or\or \hskip 1em \or \hskip 2em \else \hskip 3em \fi%
      #6\nobreak\relax
    \dotfill\hbox to\@pnumwidth{\@tocpagenum{#7}}\par
    \nobreak
    \endgroup
  \fi}
\makeatother

%redefine paragraph to subsubsubsection level
\makeatletter
\expandafter\newcommand\csname r@tocindent4\endcsname{4in} %need tocline for paragraph is it's going to be included in toc
\def\paragraph{\@startsection{paragraph}{4}%redefine paragraph heading to be bolded
  \z@\z@{-\fontdimen2\font}%
  {\normalfont\bfseries}}
  \def\l@paragraph{\@tocline{4}{0pt}{1pc}{9pc}{}} %redefine paragraph to be below subsubsection 
\makeatother

%section numbering (down to paragraph level)
\setcounter{secnumdepth}{4}
\renewcommand{\thesubsection}{\alph{subsection}}
\renewcommand{\thesubsubsection}{\roman{subsubsection}}
\renewcommand{\theparagraph}{\arabic{paragraph}}
\makeatletter
\renewcommand{\p@subsection}{\thesection.}
\renewcommand{\p@subsubsection}{\thesection.\thesubsection.}
\renewcommand{\p@paragraph}{\thesection.\thesubsection.\thesubsubsection.}
\makeatother

%create invisible section and subsection (for including pdf subdocs)
\newcommand\invisiblesection[1]{%
  \refstepcounter{section}%
  \refstepcounter{page}  %%only if numbering pages within sections
  \addcontentsline{toc}{section}{\protect\numberline{\thesection.}#1}%
  \sectionmark{#1}}

\newcommand\invisiblesubsection[1]{%
  \refstepcounter{subsection}%
  \addcontentsline{toc}{subsection}{\thesubsection.{\hskip 1em}#1}%
  %\sectionmark{#1}}
 }
 
%page numbering by section
\numberwithin{page}{section}
\renewcommand{\thepage}{\thesection-\arabic{page}}

%% Make sure that page starts from 1 with every \section
\usepackage{etoolbox}
\makeatletter
\patchcmd{\@sect}% <cmd>
  {\protected@edef}% <search>
  {\def\arg{#1}\def\arg@{section}%
   \ifx\arg\arg@\stepcounter{page}\fi%
   \protected@edef}% <replace>
  {}{}% <success><failure>
\makeatother

%hyperref load has to come after section  patching
\usepackage[breaklinks=true,colorlinks=true,citecolor=red]{hyperref}
\newcommand{\nlhref}[1]{\href{#1}{\nolinkurl{#1}}} %automatically create exact url with href

%create boolean for toggling inclusion of backmatter
\providetoggle{backmatter}
\settoggle{backmatter}{true} %set to true to toggle on/false off

%bibliography handling - has to be numeric for anonymized proposals!
\usepackage[citestyle=numeric-comp,bibstyle=numeric, sorting=none, natbib=true, backend=bibtex,hyperref=true,maxbibnames=99,maxcitenames=2]{biblatex}
\addbibresource{refs.bib, Main.bib, local.bib}


% private defs
\def\mf{\mathbf}
\def\mb{\mathbb}
\def\mc{\mathcal}
\newcommand{\mfbar}[1]{\mf{\bar{#1}}}
\newcommand{\mfhat}[1]{\mf{\hat{#1}}}
\newcommand{\refeq}[1]{Equation  (\ref{#1})} 
\newcommand{\refsec}[1]{\S\ref{#1}} 
\newcommand{\reftable}[1]{Table \ref{#1}} 
\newcommand{\refch}[1]{Chapter  \ref{#1}} 
\newcommand{\reffig}[1]{Figure \ref{#1}}
\newcommand{\intd}[1]{\ensuremath{\,\mathrm{d}#1}}
\newcommand{\leftexp}[2]{{\vphantom{#2}}^{#1}\!{#2}}
\newcommand{\leftsub}[2]{{\vphantom{#2}}_{#1}\!{#2}}
\newcommand{\fddt}[1]{\ensuremath{\leftexp{\mathcal{#1}}{\frac{\mathrm{d}}{\mathrm{d}t}}}}
\newcommand{\fdddt}[1]{\ensuremath{\leftexp{\mathcal{#1}}{\frac{\mathrm{d}^2}{\mathrm{d}t^2}}}}

%normal spacing
\renewcommand{\baselinestretch}{1}\normalsize

\begin{document}

\section{Proposal Summary}

\clearpage

%!!!EDIT
\begin{center}
\textsc{\textbf{Title}\\
\bigskip
Solicitation \#/Name}
\end{center}

\section{Table of Contents}
\renewcommand\contentsname{}
\tableofcontents

\clearpage


\section{Scientific/Technical/Management}

\subsection{Objectives and Expected Significance}\label{sec:obj}

\subsection{Perceived Impact}\label{sec:impact}

\subsection{Relevance of Proposed Work}\label{sec:relevance}

\subsection{Technical Approach and Methodology} \label{sec:tech}

\subsubsection{Background}\label{sec:back}  

\paragraph{Deeper Background}

\subsection{Potential Risks and Mitigation Strategies}\label{sec:risks}


\subsection{Plan of Work and Management Structure}\label{sec:plan}

%milestones table (optional)
 {
\renewcommand{\arraystretch}{1.1}
\begin{longtable}{|m{0.03\textwidth}|m{0.65\textwidth}|m{0.25\textwidth}|}
\caption{Major project milestones.\label{tbl:milestones}}\\
\hline
 & \textbf{Milestone} & \textbf{Expected} \\
\hline
\endfirsthead
\multicolumn{3}{c}%
{\tablename\ \thetable\ -- \textit{Continued from previous page}} \\
\hline
 & \textbf{Milestone} & \textbf{Expected} \\
\hline
\endhead
\hline \multicolumn{3}{r}{\textit{Continued on next page}} \\
\endfoot
\hline
\endlastfoot
1 & Win Award & 2020\\
\hline
2 & Get Money & 2025\\
\hline
3& Complete Work &  2028\\
\hline
\end{longtable}
}

%NB: some ROSES elements want this here, some want a separate DMP
\subsection{Data Management Plan}\label{sec:dmp}
\subsubsection{Data Types, Volume, and Formats}

\subsubsection{Data Archiving}

\subsubsection{Mechanisms for Long-term Preservation and Public Access}

\subsubsection{Schedule for Data Archiving and Sharing}  


\clearpage
\section{References}
\printbibliography[heading=none]

%%most of the backmatter will be moved to unredacted separate doc
\iftoggle{backmatter}{

\clearpage
\section{Personnel and Work Effort Table} %fully anonymized
\renewcommand*{\arraystretch}{1.5}
\noindent\begin{tabular}{| L{1.35in} | L{1.65in} | L{1.5in} | L{1.2in} |}
\hline
\textbf{Person or Role} &
\textbf{Time Charged to this Proposal} & 
\textbf{Time not charged to this proposal} & 
\textbf{Total time per person/year} \\
\hline
PI (Institution 1) & 0.5 academic months/year & N/A & 0.5 months/year \\
\hline
GRA 1 (Institution 1) & 12 months/year & N/A & 12 months/year\\
\hline
Co-I 1 (Institution 2) & 1 month/year & N/A & 1 month/year\\
\hline
GRA 2 (Institution 2) & 12 months/year & N/A & 12 months/year\\
\hline
Co-I 2   (Institution 3) & 1 months/year & N/A & 1 months/year  \\
\hline
GRA 3 (Institution 3) & 3 months/year & N/A & 3 months/year\\
\hline
Collaborator 1 &  N/A  & de minimis & de minimis \\
\hline
\end{tabular}

\clearpage
%remember - these must all be anonymous AND redacted for salary/fringe
\invisiblesection{Budget Narrative}
\invisiblesubsection{Institution 1 (lead)}
%\includepdf[pages=-,pagecommand={}]{inst1_budget_justification_redacted}
\invisiblesubsection{Institution 2}
%\includepdf[pages=-,pagecommand={}]{inst2_budget_justification_redacted}
\invisiblesubsection{Institution 3}
%\includepdf[pages=-,pagecommand={}]{inst3_budget_justification_redacted}

\clearpage
\invisiblesection{Detailed Budget}
\invisiblesubsection{Institution 1 (lead)}
%\includepdf[pages=-,pagecommand={}]{inst1_detailed_budget_redacted}
\invisiblesubsection{Institution 2}
%\includepdf[pages=-,pagecommand={}]{inst2_detailed_budget_redacted}
\invisiblesubsection{Institution 3}
%\includepdf[pages=-,pagecommand={}]{inst3_detailed_budget_redacted}
}


%Separate PDF documents from proposal document
%"Total" Budget Document (separate PDF file attached as document type "Total Budget")
%"Non-anoynmized expertise and resources doc


\end{document}